\chapter{ESTADO DEL ARTE}

\section{Estado del Arte}

El estado del arte del proyecto ITPA: Monitorizador IOT combina conocimientos de agricultura, tecnologías IoT (Internet de las Cosas) y sistemas de iluminación artificial, proporcionando un contexto para la implementación de soluciones avanzadas en invernaderos. A continuación, se describen las áreas clave:

\subsection{Iluminación Artificial y Crecimiento Vegetal}
Numerosos estudios han explorado cómo la luz artificial influye en el crecimiento de las plantas. Investigaciones como las de González y Hernández (2019) y Zhang y Liu (2021) destacan el impacto positivo de la manipulación del espectro lumínico para optimizar el desarrollo de diferentes especies vegetales. Los sistemas basados en luces LED han mostrado ser particularmente efectivos debido a su capacidad para ajustar la intensidad y el espectro según las necesidades específicas de las plantas. Sin embargo, estos estudios han sido limitados a entornos experimentales y requieren soluciones escalables para su aplicación en sistemas agrícolas más amplios.

\subsection{IoT en la Agricultura}
El Internet de las Cosas ha revolucionado la agricultura al permitir la monitorización y control en tiempo real de variables ambientales críticas. Proyectos como los de Xiang et al. (2021) han demostrado cómo los sensores integrados en invernaderos pueden ajustar automáticamente las condiciones ambientales basándose en datos en tiempo real. Aunque estas tecnologías han optimizado la productividad agrícola, persisten desafíos en la integración eficiente de dispositivos IoT para la toma de decisiones basadas en análisis predictivos y gráficos históricos.

\subsection{Frameworks IoT para Agricultura}
Los frameworks actuales para IoT en la agricultura están diseñados para ofrecer soluciones modulares y escalables. Fernández (2021) y López y Torres (2020) han desarrollado plataformas que permiten la recolección, procesamiento y análisis de datos de sensores, pero carecen de herramientas específicas para la personalización de modelos y la generación automática de gráficos históricos y descargables. Estas funcionalidades son esenciales para mejorar la gestión de invernaderos y optimizar el uso de recursos.

\subsection{Proyecto ITPA: Monitorizador IOT y su Propuesta de Valor}
El proyecto ITPA: Monitorizador IOT del Tecnológico de Pabellón de Arteaga aborda estas brechas mediante el desarrollo de un framework IoT que combina tecnologías modernas como Node.js, TypeScript y MongoDB con sistemas de iluminación artificial avanzados. Su enfoque incluye:
\begin{itemize}
    \item \textbf{Integración de IoT:} Monitoreo en tiempo real de variables como temperatura, humedad y niveles de luz, utilizando sensores y dispositivos como Raspberry Pi.
    \item \textbf{Análisis Avanzado:} Capacidad para graficar, agrupar por fechas y descargar datos en diversos formatos.
    \item \textbf{Escalabilidad y Flexibilidad:} Un sistema modular que puede adaptarse a las necesidades específicas de distintas especies vegetales y a otros entornos industriales.
\end{itemize}

Con esta propuesta, ITPA: Monitorizador IOT no solo mejora la productividad agrícola en entornos controlados, sino que también fomenta la sostenibilidad al optimizar el uso de recursos y reducir el impacto ambiental. Esto coloca al proyecto en la vanguardia de la agricultura moderna y sostenible.
